\chapter{بازهدایت جریان‌های ورودی و خروجی (I/O Redirection)}

در این بخش، به بررسی یکی از کلیدی‌ترین و قدرتمندترین قابلیت‌های رابط خط فرمان، یعنی بازهدایت ورودی/خروجی (\en{I/O Redirection}) می‌پردازیم. این سازوکار به شما اجازه می‌دهد تا جریان داده‌های ورودی و خروجی دستورات را مدیریت کرده و آن‌ها را به جای ارسال به نمایشگر یا دریافت از صفحه‌کلید، به سمت فایل‌ها یا سایر فرامین هدایت کنید.

در این فصل با ابزارهای تخصصی زیر آشنا خواهیم شد:

\begin{itemize}
  \item \cmd{cat}: جهت الحاق و نمایش محتوای فایل‌ها.
  \item \cmd{sort}: برای مرتب‌سازی خطوط متنی بر اساس معیارهای مختلف.
  \item \cmd{uniq}: جهت شناسایی، گزارش یا حذف خطوط تکراری در یک فایل.
  \item \cmd{grep}: برای جستجو و استخراج خطوط منطبق بر یک الگوی مشخص.
  \item \cmd{wc}: جهت محاسبه و شمارش تعداد خطوط، واژگان و بایت‌ها.
  \item \cmd{head}: برای مشاهده و بررسی بخش‌های ابتدایی یک فایل.
  \item \cmd{tail}: جهت پایش و نمایش بخش‌های انتهایی یک فایل.
  \item \cmd{tee}: جهت خواندن از ورودی استاندارد و نوشتن هم‌زمان در خروجی استاندارد و فایل‌ها.
\end{itemize}

\section{جریان‌های استاندارد ورودی، خروجی و خطا (Standard Streams)}

بسیاری از ابزارهای خط فرمان، خروجی خود را به دو جریان مجزا تفکیک می‌کنند:

\begin{itemize}
  \item \textbf{نتایج عملیاتی برنامه:} داده‌های اصلی و محوری که ابزار مذکور با هدف تولید یا پردازش آن‌ها طراحی شده است.
  \item \textbf{پیام‌های وضعیت و خطا:} گزارش‌هایی که اطلاعاتی درباره‌ی چگونگی اجرای فرآیند، هشدارها یا اختلالات احتمالی در اختیار کاربر قرار می‌دهند.
\end{itemize}

مطابق با فلسفه‌ی یونیکس که بر اصل «همه چیز یک فایل است» استوار است، این جریان‌ها به‌جای ارسال مستقیم به سخت‌افزار نمایشگر، به فایل‌های انتزاعی ویژه‌ای تحت عنوان خروجی استاندارد (\en{stdout}) و خطای استاندارد (\en{stderr}) هدایت می‌شوند. به‌صورت پیش‌فرض، هر دوی این جریان‌ها به محیط کنسول یا ترمینال شما متصل شده‌اند.

علاوه بر این، اکثر برنامه‌ها داده‌های ورودی خود را از منبعی موسوم به ورودی استاندارد (\en{stdin}) دریافت می‌کنند که در حالت پیش‌فرض به صفحه‌کلید متصل است. بازهدایت ورودی/خروجی (\en{I/O Redirection}) سازوکاری است که امکان بازتعریف این مسیرهای پیش‌فرض را فراهم می‌آورد؛ بدین معنا که می‌توان تعیین کرد ورودی از چه منبعی خوانده شده و خروجی یا خطا به چه مقصدی فرستاده شود.

\section{بازهدایت خروجی استاندارد (stdout)}

قابلیت بازهدایت به ما این امکان را می‌دهد تا مقصد خروجی استاندارد (\en{stdout}) را از محیط نمایشگر به یک فایل یا هر مقصد دیگری که سیستم‌عامل آن را به‌عنوان یک فایل می‌شناسد تغییر دهیم. برای تحقق این امر، از عملگر بازهدایت \cmd{>} به‌همراه نام فایل مقصد استفاده می‌شود.

ذخیره‌سازی خروجیِ یک فرمان در قالب یک فایل، از جنبه‌های مختلفی نظیر مستندسازی یا پردازش‌های آتی بسیار سودمند است. به‌طور نمونه، جهت انتقال و ثبت نتایج حاصل از اجرای دستور \cmd{ls -l /usr/bin} در فایلی تحت عنوان \file{ls-output.txt}، فرمان زیر را اجرا می‌کنیم:

\begin{latin}
\begin{lstlisting}
[me@linuxbox ~]$ ls -l /usr/bin > ls-output.txt
\end{lstlisting}
\end{latin}

در این سناریو، هیچ داده‌ای بر روی نمایشگر چاپ نمی‌شود و تمامی نتایج به فایل \file{ls-output.txt} منتقل می‌گردد. برای حصول اطمینان از ایجاد فایل و بررسی حجم داده‌های انتقال یافته، می‌توان از دستور \cmd{ls -l} استفاده کرد:

\begin{latin}
\begin{lstlisting}
[me@linuxbox ~]$ ls -l ls-output.txt
-rw-rw-r-- 1 me       me          167878 2025-02-01 15:07 ls-output.txt
\end{lstlisting}
\end{latin}

همچنین جهت بازبینی و مطالعه‌ی محتویات ثبت شده، ابزار \cmd{less} گزینه‌ی مناسبی خواهد بود:

\begin{latin}
\begin{lstlisting}
[me@linuxbox ~]$ less ls-output.txt
\end{lstlisting}
\end{latin}

درک این نکته حیاتی است که عملگر بازهدایت \cmd{>} صرفاً خروجی استاندارد (\en{stdout}) را بازهدایت می‌کند؛ در حالی که پیام‌های سیستمی و خطا به مجرای خطای استاندارد (\en{stderr}) ارسال می‌شوند. به همین سبب، چنانچه دستوری را با خطای ساختاری اجرا کنید، پیام خطا همچنان بر روی نمایشگر ظاهر شده و به فایل منتقل نمی‌گردد:

\begin{latin}
\begin{lstlisting}
[me@linuxbox ~]$ ls -l /bin/usr > ls-output.txt
ls: cannot access /bin/usr: No such file or directory
\end{lstlisting}
\end{latin}

در این سناریو، فایل \file{ls-output.txt} تهی باقی می‌ماند؛ زیرا دستور \cmd{ls} هیچ خروجی استانداردی (\en{stdout}) تولید نکرده و صرفاً یک گزارش خطا صادر نموده است. این رفتار منطبق بر استانداردهای یونیکس است که در آن، فرامین پیام‌های خطای خود را به جریانی مجزا از نتایج اصلی می‌فرستند تا کاربر همواره از وقوع اختلالات مطلع باشد.

اکنون وضعیت فایل مقصد را بررسی می‌کنیم:

\begin{latin}
\begin{lstlisting}
[me@linuxbox ~]$ ls -l ls-output.txt
-rw-rw-r-- 1 me       me            0 2025-02-01 15:08 ls-output.txt
\end{lstlisting}
\end{latin}

همان‌طور که ملاحظه می‌شود، حجم فایل به صفر بایت کاهش یافته است. علت این پدیده آن است که هنگام بهره‌گیری از عملگر بازهدایت (\cmd{>})، سیستم پیش از اجرای دستور، فایل مقصد را باز کرده و محتوای آن را جهت نوشتن داده‌های جدید پاک‌سازی می‌کند. در این مثال، از آنجا که دستور با خطا مواجه شد و جریانی در خروجی استاندارد ایجاد نگردید، فرآیند بازهدایت با یک فایل خالی خاتمه یافت.

در واقع، چنانچه نیاز به پاک‌سازی کامل محتوای یک فایل موجود یا ایجاد یک فایل تهی جدید داشته باشیم، می‌توانیم از یک ترفند کارآمد در پوسته استفاده کنیم. با به‌کارگیری عملگر بازهدایت (\cmd{>}) بدون درج هیچ دستوری پیش از آن، این هدف محقق می‌شود:

\begin{latin}
\begin{lstlisting}
[me@linuxbox ~]$ > ls-output.txt
\end{lstlisting}
\end{latin}

حال این پرسش کلیدی مطرح می‌شود که چگونه می‌توان خروجی یک فرمان را به فایلی الحاق کرد، به‌گونه‌ای که داده‌های پیشین حفظ شوند؟ برای این منظور، از عملگر الحاق (\en{Appendance}) با نماد \cmd{>>} استفاده می‌کنیم:

\begin{latin}
\begin{lstlisting}
[me@linuxbox ~]$ ls -l /usr/bin >> ls-output.txt
\end{lstlisting}
\end{latin}

استفاده از عملگر الحاق \cmd{>>} موجب می‌شود که خروجی استاندارد به‌جای بازنویسی (\en{Overwrite})، به انتهای فایل مقصد افزوده شود. لازم به ذکر است که اگر فایل هدف از پیش وجود نداشته باشد، این عملگر همانند عملگر بازهدایت \cmd{>} اقدام به ساخت آن می‌کند. جهت اعتبارسنجی این مفهوم، فرمان فوق را چندین بار متوالی تکرار می‌کنیم:

\begin{latin}
\begin{lstlisting}
[me@linuxbox ~]$ ls -l /usr/bin >> ls-output.txt
[me@linuxbox ~]$ ls -l /usr/bin >> ls-output.txt
[me@linuxbox ~]$ ls -l /usr/bin >> ls-output.txt
[me@linuxbox ~]$ ls -l ls-output.txt
-rw-rw-r-- 1 me       me          503634 2025-02-01 15:45 ls-output.txt
\end{lstlisting}
\end{latin}

اکنون با بررسی مشخصات فایل، مشاهده می‌کنیم که حجم آن به‌میزان قابل‌توجهی افزایش یافته است. با پایش حجم نهایی، تأیید می‌شود که خروجی هر سه مرحله به‌صورت پیاپی در فایل ثبت شده و یکپارچگی داده‌های قبلی حفظ گشته است.

\section{گروه‌بندی دستورات (Group Commands)}

با تکیه بر آموخته‌های پیشین، جهت اجرای مجموعه‌ای از دستورات و هدایت خروجی آن‌ها به یک فایل گزارش (\en{Log File})، می‌توان از رویکرد زیر استفاده کرد:

\begin{latin}
\begin{lstlisting}
[me@linuxbox ~]$ command1 > logfile.txt
[me@linuxbox ~]$ command2 >> logfile.txt
[me@linuxbox ~]$ command3 >> logfile.txt
\end{lstlisting}
\end{latin}

در این روش، دستور نخست فایل \file{logfile.txt} را ایجاد یا بازنویسی (\en{Overwrite}) کرده و دستورات متعاقب، نتایج خود را به انتهای آن الحاق (\en{Append}) می‌کنند. اگرچه این شیوه عملیاتی است، اما به‌دلیل تکرار فرآیند نگارش و احتمال بروز خطای انسانی، بهینه به‌نظر نمی‌رسد. قطعاً راهکار هوشمندانه‌تری در محیط پوسته پیش‌بینی شده است.

همان‌طور که پیش‌تر اشاره شد، امکان اجرای متوالی چندین فرمان در یک خط واحد با بهره‌گیری از جداکننده‌ی سمی‌کالن (\cmd{;}) میسر است:

\begin{latin}
\begin{lstlisting}
[me@linuxbox ~]$ command1; command2; command3
\end{lstlisting}
\end{latin}

بر همین اساس، می‌توان توالی فرامین و بازهدایت‌های مربوطه را نیز در یک سطر ادغام کرد:

\begin{latin}
\begin{lstlisting}
[me@linuxbox ~]$ command1 > logfile.txt; command2 >> logfile.txt; command3 >> logfile.txt
\end{lstlisting}
\end{latin}

با این حال، چنانچه بخواهیم کل این سلسله‌مراتب را به‌عنوان یک واحد منطقی با یک جریان خروجی یکپارچه تلقی کنیم، بهترین راهکار استفاده از «گروه‌بندی فرامین» است. در این روش، تمامی فرامین درون آکولاد \cmd{\{ \}} محصور می‌شوند:

\begin{latin}
\begin{lstlisting}
[me@linuxbox ~]$ { command1; command2; command3; } > logfile.txt
\end{lstlisting}
\end{latin}

با این تکنیک، پوسته کل توالی محصور در آکولاد را به‌عنوان یک موجودیت واحد پردازش کرده و خروجی تجمیعی آن‌ها را به‌صورت یکجا به فایل \file{logfile.txt} هدایت می‌کند. شایان ذکر است که رعایت دو نکته‌ی نحوی (\en{Syntax}) برای اجرای صحیح الزامی است: نخست، درج فاصله‌ی خالی (\en{Space}) پس از آکولاد باز و پیش از آکولاد بسته، و دوم، خاتمه یافتن آخرین دستور درون گروه با استفاده از نویسه سمی‌کالن (\cmd{;}) یا خط جدید.

\section{بازهدایت خطای استاندارد (stderr)}

بازهدایت جریان خطای استاندارد (\en{stderr}) نسبت به خروجی استاندارد پیچیدگی بیشتری دارد؛ زیرا مستلزم ارجاع مستقیم به توصیف‌گر فایل (\en{File Descriptor}) است. در محیط سیستم‌عامل، برنامه‌ها می‌توانند داده‌های خود را به چندین جریان مجزا که با شناسه‌های عددی متمایز می‌شوند، ارسال کنند. سه جریان بنیادین در لینوکس عبارتند از:

\begin{itemize}
  \item \cmd{0}: ورودی استاندارد (\en{stdin})
  \item \cmd{1}: خروجی استاندارد (\en{stdout})
  \item \cmd{2}: خطای استاندارد (\en{stderr})
\end{itemize}

پوسته (\en{Shell}) از نحوِ (\en{Syntax}) خاصی برای هدایت این جریان‌ها بر اساس شماره‌ی توصیف‌گر آن‌ها استفاده می‌کند. از آنجا که جریان خطا با شناسه \cmd{2} شناخته می‌شود، می‌توان عملیات بازهدایت را به شکل زیر انجام داد:

\begin{latin}
\begin{lstlisting}
[me@linuxbox ~]$ ls -l /bin/usr 2> ls-error.txt
\end{lstlisting}
\end{latin}

در این فرمان، عدد \cmd{2} دقیقاً پیش از عملگر بازهدایت \cmd{>} درج می‌شود تا به پوسته ابلاغ کند که صرفاً جریان خطای استاندارد باید به فایل \file{ls-error.txt} منتقل شود. با این رویکرد، پیام‌های خطا به‌جای چاپ بر روی نمایشگر، در فایل مذکور بایگانی می‌شوند.

\section{بازهدایت هم‌زمان خروجی استاندارد و خطای استاندارد}

برای ثبت تمامی نتایج حاصل از اجرای یک فرمان در یک فایل منفرد، لازم است هر دو جریان خروجی استاندارد (\en{stdout}) و خطای استاندارد (\en{stderr}) به‌طور هم‌زمان بازهدایت شوند. برای تحقق این هدف، دو رویکرد اصلی وجود دارد.

\textbf{روش کلاسیک (سازگار با نسخه‌های قدیمی پوسته):} در این روش، دو عملیات بازهدایت به‌صورت متوالی صورت می‌پذیرد. ابتدا جریان خروجی استاندارد (\en{stdout}) به فایل مقصد هدایت شده و در گام بعد، جریان خطای استاندارد (\en{stderr}) به سمت خروجی استاندارد (\en{stdout}) بازتعریف می‌شود. این فرآیند با استفاده از عبارت ترکیبی \cmd{2>\&1} انجام می‌گیرد:

\begin{latin}
\begin{lstlisting}
[me@linuxbox ~]$ ls -l /bin/usr > ls-output.txt 2>&1
\end{lstlisting}
\end{latin}

در این دستور، عملگر بازهدایت \cmd{>} ابتدا خروجی استاندارد را به فایل \file{ls-output.txt} منتقل می‌کند. سپس عبارت \cmd{2>\&1} به پوسته فرمان می‌دهد که جریان خطای استاندارد (توصیف‌گر فایل \cmd{2}) را به همان مقصدی هدایت کند که خروجی استاندارد (توصیف‌گر فایل \cmd{1}) در حال ارسال به آن است.

شایان ذکر است که ترتیب قرارگیری این عبارات در خط فرمان بسیار حیاتی است. چنانچه این ترتیب جابه‌جا شود، جریان خطای استاندارد (\en{stderr}) به‌جای فایل، به نمایشگر هدایت خواهد شد؛ زیرا در لحظه‌ی ارجاع، جریان خروجی استاندارد (\en{stdout}) هنوز به فایل متصل نشده و همچنان به نمایشگر (مبدأ پیش‌فرض خود) اشاره دارد.

\textbf{روش مدرن (مختص نسخه‌های جدید \en{Bash}):} نسخه‌های جدیدتر پوسته \en{Bash}، نحوِ (\en{Syntax}) ساده‌تری را برای تسهیل این فرآیند معرفی کرده‌اند. با بهره‌گیری از نماد ترکیبی \cmd{\&>}، می‌توانید هر دو جریان خروجی استاندارد و خطای استاندارد را به‌طور هم‌زمان به یک فایل واحد هدایت کنید:

\begin{latin}
\begin{lstlisting}
[me@linuxbox ~]$ ls -l /bin/usr &> ls-output.txt
\end{lstlisting}
\end{latin}

این دستور از نظر عملکردی دقیقاً معادل روش کلاسیک است و تمامی خروجی‌ها (اعم از نتایج و خطاها) را در فایل \file{ls-output.txt} تجمیع می‌کند. علاوه بر این، چنانچه قصد داشته باشید هر دو جریان را به‌جای بازنویسی، به انتهای یک فایل موجود اضافه (\en{Append}) کنید، می‌توانید از نماد \cmd{\&>>} استفاده نمایید:

\begin{latin}
\begin{lstlisting}
[me@linuxbox ~]$ ls -l /bin/usr &>> ls-output.txt
\end{lstlisting}
\end{latin}

این رویکرد مدرن، فرآیند ثبت و بایگانی تمام خروجی‌های عملیاتی را به شکلی فشرده، خوانا و با ضریب خطای کمتر امکان‌پذیر می‌سازد.

\section{مدیریت خروجی‌های بلااستفاده و امحای داده‌ها}

در بسیاری از سناریوها، خروجی‌های حاصل از یک فرمان، به‌ویژه پیام‌های وضعیت یا گزارش‌های خطا، فاقد ارزش عملیاتی هستند و صرفاً موجب شلوغی محیط ترمینال می‌شوند. در چنین مواردی، می‌توان از قابلیتی جهت پنهان‌سازی و امحای این داده‌ها استفاده کرد. برای این منظور، خروجی مورد نظر به فایل ویژه‌ای تحت عنوان \file{/dev/null} هدایت می‌شود.

فایل \file{/dev/null} یک دستگاه مجازی در سلسله‌مراتب سیستم‌عامل است که به‌عنوان یک «سیاه‌چاله داده» عمل می‌کند؛ به این معنا که هر جریانی به آن ارسال شود، بلافاصله و بدون اشتغال فضای دیسک یا پردازش ثانویه، از بین می‌رود. در ادبیات فنی لینوکس، این فایل را اغلب \en{``bit bucket''} یا سطل زباله بیت‌ها می‌نامند.

به‌عنوان مثال، جهت جلوگیری از نمایش پیام‌های خطای ناشی از دستور \cmd{ls} و پاک‌سازی محیط نمایشگر، جریان خطا را به این مسیر هدایت می‌کنیم:

\begin{latin}
\begin{lstlisting}
[me@linuxbox ~]$ ls -l /bin/usr 2> /dev/null
\end{lstlisting}
\end{latin}

در این فرمان، عملگر \cmd{2>} جریان خطای استاندارد (\en{stderr}) را به سمت \file{/dev/null} منحرف می‌کند. در نتیجه، چنانچه دایرکتوری مورد نظر یافت نشود، پیام خطا به‌جای چاپ در کنسول، به‌صورت بی‌صدا امحا می‌گردد.

\section{تحلیل مفهوم /dev/null در فرهنگ یونیکس}

مفهوم \en{Bit Bucket} (سطل زباله‌ی بیت‌ها) یکی از استعاره‌های دیرین و بنیادین در فرهنگ سیستم‌عامل‌های خانواده‌ی یونیکس است. از آنجا که فایل \file{/dev/null} به‌طور گسترده برای امحای داده‌های بلااستفاده به‌کار می‌رود، فراتر از یک ابزار سیستمی، به ادبیات روزمره‌ی متخصصان راه یافته است. زمانی که یک کاربر یا توسعه‌دهنده به مزاح یا کنایه بیان می‌کند که «نظرات شما را به \file{/dev/null} می‌فرستد»، در واقع به نادیده گرفتن کامل و عدم توجه به آن پیام‌ها اشاره دارد. این اصطلاح به‌خوبی جایگاه \file{/dev/null} را به‌عنوان نمادی برای حذفِ مطلق، بی‌صدا و بازیابی‌ناپذیر در دنیای پردازش داده‌ها تبیین می‌کند.

\section{بازهدایت ورودی استاندارد (stdin)}

تا بدین‌جا، تمرکز ما عمدتاً بر مدیریت و بازهدایت جریان‌های خروجی معطوف بود. اکنون زمان آن فرا رسیده است تا به بررسی ماهیت و نحوه تعامل با ورودی استاندارد (\en{stdin}) بپردازیم. اگرچه ممکن است پیش از این به‌طور غیرمستقیم با فرامینی که از ورودی استاندارد (\en{stdin}) بهره می‌برند کار کرده باشید، اما برای تبیین دقیق این مفهوم، از یکی از بنیادی‌ترین ابزارهای لینوکس یعنی دستور \cmd{cat} استفاده می‌کنیم.

در ساختار عملیاتی لینوکس، همان‌طور که می‌توان خروجی را از نمایشگر به سمت یک فایل منحرف کرد، این امکان نیز وجود دارد که منبع تأمین داده برای یک برنامه را از صفحه‌کلید به یک فایل متنی تغییر داد. این فرآیند زیربنای خودکارسازی بسیاری از وظایف سیستمی است.

\section{دستور cat: الحاق و نمایش محتوای فایل‌ها}

دستور \cmd{cat} که مخفف عبارت \en{catenate} (به معنای الحاق کردن) است، ابزاری کلیدی برای خواندن محتوای یک یا چند فایل و ارسال مستقیم آن‌ها به خروجی استاندارد (\en{stdout}) محسوب می‌شود. نحو این فرمان به صورت زیر است:

\begin{latin}
\begin{lstlisting}
cat [file...]
\end{lstlisting}
\end{latin}

این ابزار به‌طور گسترده جهت بازبینی سریع محتوای فایل‌های متنی کم‌حجم که نیازی به صفحه‌بندی (\en{Paging}) ندارند، استفاده می‌شود. برای نمونه، جهت مشاهده‌ی مندرجات فایل \file{ls-output.txt} از فرمان زیر استفاده می‌کنیم:

\begin{latin}
\begin{lstlisting}
[me@linuxbox ~]$ cat ls-output.txt
\end{lstlisting}
\end{latin}

از آنجا که \cmd{cat} قادر است چندین فایل را به‌عنوان آرگومان ورودی بپذیرد، می‌توان از آن برای ادغام محتوای چندین سند در قالب یک جریان داده‌ی واحد بهره برد. این ویژگی، \cmd{cat} را به ابزاری بسیار منعطف و کارآمد در فرآیندهای بازهدایت ورودی و خروجی مبدل کرده است.

تصور کنید فایلی حجیم را دانلود کرده‌اید که به چندین بخش مجزا تقسیم شده است؛ مشابه رویه‌ای که در توزیع فایل‌های چندرسانه‌ای در بسترهایی مانند یوزنت (\en{Usenet}) رایج است. برای بازسازی و اتصال مجدد این قطعات، می‌توانید از یک راهکار هوشمندانه مبتنی بر دستور \cmd{cat} و نویسه‌های جانشین (\en{Wildcards}) بهره بگیرید.

چنانچه قطعات فایل با الگوی نام‌گذاری \file{movie.mpeg.001} تا \file{movie.mpeg.099} مشخص شده باشند، با اجرای فرمان زیر می‌توان آن‌ها را یکپارچه کرد:

\begin{latin}
\begin{lstlisting}
cat movie.mpeg.0* > movie.mpeg
\end{lstlisting}
\end{latin}

این دستور از یکی از قابلیت‌های کلیدی پوسته (\en{Shell}) استفاده می‌کند: پیش از اجرای قطعی فرمان، پوسته نویسه‌ی جانشین را با تمامی نام‌فایل‌های منطبق جایگزین می‌کند. نکته‌ی حائز اهمیت این است که این جایگزینی همواره بر اساس ترتیب الفبایی/عددی صورت می‌پذیرد. در نتیجه، فایل‌ها با توالی صحیح از \file{001} تا \file{099} به عنوان آرگومان به دستور \cmd{cat} تحویل داده می‌شوند. سپس \cmd{cat} محتوای آن‌ها را به ترتیبِ دریافت، خوانده و جریان خروجی یکپارچه را به سمت فایل مقصد (\file{movie.mpeg}) هدایت می‌کند. نتیجه‌ی این فرآیند، بازسازی کامل فایل اصلی خواهد بود.

اگرچه این مکانیزم کاملاً منطقی به نظر می‌رسد، اما ممکن است این پرسش مطرح شود که این فرآیند چه پیوندی با مفهوم ورودی استاندارد (\en{stdin}) دارد؟ تا این مرحله، ارتباط مستقیمی مشهود نیست؛ اما بیایید برای درک عمیق‌تر این موضوع، یک آزمایش عملی انجام دهیم.

چنانچه دستور \cmd{cat} را بدون ارجاع به هیچ فایل خاصی اجرا کنید، به نظر می‌رسد که برنامه در وضعیت تعلیق قرار گرفته است:

\begin{latin}
\begin{lstlisting}
[me@linuxbox ~]$ cat

\end{lstlisting}
\end{latin}

این رفتار، در واقع عملکرد استاندارد و پیش‌بینی‌شده‌ی برنامه است. هنگامی که \cmd{cat} بدون آرگومان فراخوانی می‌شود، منبع داده‌ی خود را از فایل به ورودی استاندارد (\en{stdin}) تغییر می‌دهد. از آنجا که این جریان به‌طور پیش‌فرض به صفحه‌کلید متصل است، برنامه در انتظار دریافت ورودی از سوی کاربر باقی می‌ماند. متن دلخواه خود را تایپ کرده و کلید \en{Enter} را فشار دهید:

\begin{latin}
\begin{lstlisting}
[me@linuxbox ~]$ cat
The quick brown fox jumps over the lazy dog.
\end{lstlisting}
\end{latin}

اکنون برای اعلام خاتمه‌ی فایل (\en{End-Of-File} یا به‌اختصار \en{EOF})، ترکیب کلیدی \cmd{Ctrl+d} را بفشارید. در این حالت، \cmd{cat} داده‌های دریافتی از ورودی استاندارد را به خروجی استاندارد (\en{stdout}) کپی می‌کند؛ در نتیجه، متنی که تایپ کرده‌اید، مجدداً بر روی نمایشگر ظاهر می‌شود:

\begin{latin}
\begin{lstlisting}
[me@linuxbox ~]$ cat
The quick brown fox jumps over the lazy dog.
The quick brown fox jumps over the lazy dog.
\end{lstlisting}
\end{latin}

در غیاب آرگومانِ نام فایل، \cmd{cat} به‌عنوان یک پل میان ورودی و خروجی استاندارد عمل کرده و محتوای دریافتی را تکرار می‌کند. از این ویژگی می‌توان برای ایجاد سریع فایل‌های متنی کوتاه بهره برد.

به‌منظور تولید یک فایل متنی با محتوای دلخواه، می‌توان دستور \cmd{cat} را با عملگر بازهدایت خروجی (\cmd{>}) ترکیب کرد. با این تکنیک، جریانی که از ورودی استاندارد (صفحه‌کلید) دریافت می‌شود، به‌جای نمایشگر، مستقیماً به سمت فایل مقصد هدایت و در آن ذخیره می‌گردد.

به منظور تولید فایلی تحت عنوان \file{lazy\_dog.txt} که حاوی محتوای مورد نظر باشد، فرمان زیر را در ترمینال اجرا کنید؛ سپس متن خود را تایپ نموده و در انتها، ترکیب کلیدی \cmd{Ctrl+d} را جهت ارسال سیگنال پایان فایل (\en{EOF}) به برنامه بفشارید:

\begin{latin}
\begin{lstlisting}
[me@linuxbox ~]$ cat > lazy_dog.txt
The quick brown fox jumps over the lazy dog.
\end{lstlisting}
\end{latin}

با این رویکرد، ما در واقع از خط فرمان به عنوان یک ویرایشگر متن (\en{Text Editor}) بسیار ساده و بدوی استفاده کرده‌ایم. جهت اعتبارسنجی و اطمینان از صحت محتویات ذخیره‌شده، می‌توان مجدداً از \cmd{cat} برای هدایت داده‌های فایل به خروجی استاندارد (نمایشگر) بهره گرفت:

\begin{latin}
\begin{lstlisting}
[me@linuxbox ~]$ cat lazy_dog.txt
The quick brown fox jumps over the lazy dog.
\end{lstlisting}
\end{latin}

در اینجا، برنامه فایل را باز کرده و محتوای آن را مستقیماً به خروجی استاندارد ارسال می‌کند.

اکنون که دریافتیم فرمان \cmd{cat} قابلیت پذیرش داده از ورودی استاندارد (\en{stdin}) را دارد، می‌توانیم منبع ورودی آن را از صفحه‌کلید به یک فایل متنی تغییر دهیم. این فرآیند با بهره‌گیری از عملگر بازهدایت ورودی (\cmd{<}) محقق می‌شود:

\begin{latin}
\begin{lstlisting}
[me@linuxbox ~]$ cat < lazy_dog.txt
The quick brown fox jumps over the lazy dog.
\end{lstlisting}
\end{latin}

در این وضعیت، با اجرای دستور \cmd{cat < lazy\_dog.txt}، پوسته (\en{Shell}) به‌جای انتظار برای دریافت ورودی از کاربر، فایل \file{lazy\_dog.txt} را به‌عنوان منبع داده به برنامه معرفی می‌کند. متعاقباً \cmd{cat} محتوای دریافتی را به خروجی استاندارد (نمایشگر) ارسال می‌نماید. همان‌طور که ملاحظه می‌کنید، نتیجه‌ی نهایی با حالتی که نام فایل را به‌عنوان آرگومان مستقیم به \cmd{cat} پاس می‌دهیم، یکسان است.

اگرچه این رویکرد در مورد دستور ساده‌ای مانند \cmd{cat} مزیت عملیاتی خاصی ندارد، اما حامل یک مفهوم بنیادین در معماری یونیکس است: هر فایلی می‌تواند به‌عنوان منبع ورودی استاندارد برای یک فرمان عمل کند. در مباحث پیش‌رو خواهیم دید که چگونه دستورات پیچیده‌تر، از این قابلیت برای پردازش‌های خودکار به شکلی بسیار کارآمدتر بهره می‌برند.

\section{سازوکار پایپ‌لاین‌ها (Pipelines)}

اوج کارآمدی فرامین در خواندن داده‌ها از ورودی استاندارد (\en{stdin}) و ارسال آن‌ها به خروجی استاندارد (\en{stdout})، در قابلیتی تحت عنوان «پایپ‌لاین» یا \en{Pipeline} تحقق می‌یابد. با بهره‌گیری از عملگر پایپ \cmd{|} (کاراکتر خط عمودی)، می‌توان خروجی استاندارد یک فرمان را مستقیماً به ورودی استاندارد فرمان بعدی متصل کرد. ساختار کلی این عملیات به شرح زیر است:

\begin{latin}
\begin{lstlisting}
command1 | command2
\end{lstlisting}
\end{latin}

برای درک عمیق‌تر این سازوکار، مثالی کاربردی را بررسی می‌کنیم. پیش‌تر با ابزار \cmd{less} آشنا شدیم که قابلیت بالایی در دریافت و پردازش داده‌ها از ورودی استاندارد دارد. با استفاده از یک پایپ‌لاین، می‌توانیم خروجی حجیم دستور \cmd{ls} را جهت مطالعه‌ی آسان‌تر و پیمایش صفحه‌به‌صفحه، به ورودی \cmd{less} هدایت کنیم:

\begin{latin}
\begin{lstlisting}
[me@linuxbox ~]$ ls -l /usr/bin | less
\end{lstlisting}
\end{latin}

این ترکیب عملیاتی بسیار بهینه است؛ چرا که امکان مشاهدهٔ صفحه‌به‌صفحه خروجی هر فرمانی را بدون نیاز به ذخیره‌سازی موقت در یک فایل واسط فراهم می‌سازد. در واقع، با این رویکرد می‌توان از \cmd{less} به‌عنوان یک ابزار نمایش متن (\en{Pager}) برای تمامی دستوراتی که خروجی متنی تولید می‌کنند، استفاده کرد.

\section{تمایز میان بازهدایت > و پایپ‌لاین |}

درک تفاوت عملکردی میان عملگر بازهدایت خروجی (\cmd{>}) و کاراکتر پایپ‌لاین (\cmd{|}) جهت مدیریت صحیح جریان داده‌ها و پیشگیری از بروز خطاهای جبران‌ناپذیر در سطح سیستم فایل، بسیار حیاتی است.

عملگر \cmd{>} خروجی استاندارد یک فرمان را به سمت یک فایل منحرف می‌کند. هدف اصلی از به‌کارگیری این عملگر، بایگانی یا ذخیره‌سازی نتایج اجرای یک دستور در قالب یک فایل متنی است.

\begin{latin}
\begin{lstlisting}
command1 > file1
\end{lstlisting}
\end{latin}

در این حالت، خروجی \cmd{command1} به جای نمایش روی صفحه، به طور کامل در \file{file1} نوشته می‌شود. نکته حیاتی این است که اگر \file{file1} از قبل وجود داشته باشد، محتوای آن به طور کامل بازنویسی خواهد شد.

عملگر پایپ‌لاین \cmd{|}، خروجی استاندارد یک فرمان را مستقیماً به عنوان ورودی استاندارد به فرمان دوم ارسال می‌کند. این سازوکار یک مجرای ارتباطی یک‌طرفه میان دو فرآیند مجزا ایجاد کرده و امکان زنجیره‌سازی دستورات را برای انجام عملیات پیچیده‌تر فراهم می‌آورد.

\begin{latin}
\begin{lstlisting}
command1 | command2
\end{lstlisting}
\end{latin}

در این ساختار، داده‌های تولید شده توسط \cmd{command1} به‌صورت آنی و در لحظه به \cmd{command2} منتقل می‌شوند تا تحت پردازش ثانویه قرار گیرند. برخلاف عملگر بازهدایت، در اینجا هیچ فایلی بر روی دیسک ایجاد نمی‌شود و تبادل داده‌ها صرفاً در فضای حافظه و بین دو برنامه صورت می‌پذیرد.

تفاوت بنیادین میان این دو عملگر در مقصد نهایی جریان داده نهفته است: عملگر \cmd{>} داده‌ها را درون یک فایل ذخیره می‌کند، در حالی که عملگر \cmd{|} آن‌ها را به عنوان ورودی به فرمانی دیگر تزریق می‌نماید. عدم تفکیک دقیق این دو مفهوم می‌تواند منجر به بروز خطاهای مهلکی در سطح سیستم شود.

برای درک عمق فاجعه، یک نمونه‌ی واقعی و آموزنده از خطای یک مدیر سیستم را بررسی می‌کنیم:

\begin{latin}
\begin{lstlisting}
# cd /usr/bin
# ls > less
\end{lstlisting}
\end{latin}

در دستور دوم، هدف کاربر احتمالاً مشاهده‌ی لیست فایل‌ها از طریق برنامه‌ی \cmd{less} بوده است، اما به اشتباه از عملگر بازهدایت استفاده کرده است. از آنجا که در مسیر \file{/usr/bin/} فایلی اجرایی به نام \cmd{less} (که در واقع خودِ برنامه‌ی \cmd{less} است) وجود دارد، عملگر \cmd{>} بدون درنگ آن فایل را با متون خروجیِ دستور \cmd{ls} جایگزین کرده است. در نتیجه، برنامه‌ی کاربردی \cmd{less} به طور کامل تخریب و با یک فایل متنی ساده جایگزین شده است.

آموزه‌ی کلیدی این سناریو چنین است: عملگر بازهدایت \cmd{>} بدون صدور هیچ‌گونه هشدار یا تاییدیه‌ای، اقدام به ایجاد یا بازنویسی فایل‌ها می‌کند. از این رو، همواره باید نسبت به مقصد جریان داده اطمینان حاصل کرد و با دقت مضاعف از این ابزار استفاده نمود تا از وقوع چنین آسیب‌های ساختاری به بدنه سیستم‌عامل جلوگیری شود.

\section{فیلتر‌ها (Filters)}

پایپ‌لاین‌ها ابزاری قدرتمند جهت اجرای عملیات پیچیده بر روی داده‌ها محسوب می‌شوند و امکان زنجیره‌سازی چندین فرمان را در یک توالی واحد فراهم می‌کنند. فرامینی که در این ساختار وظیفه‌ی تغییر، پردازش و تبدیل داده‌ها را بر عهده دارند، اصطلاحاً فیلتر (\en{Filter}) نامیده می‌شوند. یک فیلتر، داده‌ها را از ورودی استاندارد دریافت کرده، پردازش‌های لازم را روی آن‌ها اعمال نموده و سپس نتایج را به خروجی استاندارد ارسال می‌کند.

برای تبیین این مفهوم، ابزار \cmd{sort} را بررسی می‌کنیم. فرض کنید قصد داریم فهرستی جامع از تمامی برنامه‌های اجرایی موجود در دایرکتوری‌های \file{/bin} و \file{/usr/bin} تهیه کرده، آن‌ها را بر اساس حروف الفبا مرتب نماییم و در نهایت به‌صورت صفحه‌بندی‌شده بازبینی کنیم:

\begin{latin}
\begin{lstlisting}
[me@linuxbox ~]$ ls /bin /usr/bin | sort | less
\end{lstlisting}
\end{latin}

در این پایپ‌لاین، ابتدا دستور \cmd{ls} محتویات هر دو دایرکتوری را استخراج می‌کند. خروجی این دستور که شامل دو فهرست مجزا و لزوماً نامرتب است، به‌عنوان ورودی به فرمان \cmd{sort} هدایت می‌شود. ابزار \cmd{sort} تمامی داده‌های دریافتی را به‌صورت یکپارچه مرتب‌سازی کرده و جریان خروجی را به \cmd{less} تحویل می‌دهد تا امکان پیمایش آسان لیست نهایی فراهم گردد.

این مثال به‌خوبی نشان می‌دهد که چگونه \cmd{sort} در نقش یک فیلتر، جریان داده‌ها را برای دستیابی به نتیجه‌ی مطلوب مرتب‌سازی می‌کند. شایان ذکر است که فرمان \cmd{sort} ابزاری بسیار منعطف با گزینه‌های عملیاتی متعدد است که در مباحث پیشرفته‌تر (فصل ۲۰) به‌طور تفصیلی مورد بررسی قرار خواهد گرفت.

\section{دستور uniq: شناسایی و حذف سطرهای تکراری}

فرمان \cmd{uniq} به‌طور معمول در قالب یک فیلتر و در پیوند با دستور \cmd{sort} مورد استفاده قرار می‌گیرد. این ابزار، جریانی از داده‌های مرتب‌سازی‌شده را از ورودی استاندارد یا یک فایل دریافت کرده و به‌صورت پیش‌فرض، سطرهای تکراریِ متوالی را حذف می‌کند. برای تضمین اینکه فهرست نهایی فاقد موارد تکراری باشد (مثلاً برنامه‌هایی با نام مشابه که هم‌زمان در مسیرهای \file{/bin} و \file{/usr/bin} قرار دارند)، \cmd{uniq} را به انتهای پایپ‌لاین اضافه می‌کنیم:

\begin{latin}
\begin{lstlisting}
[me@linuxbox ~]$ ls /bin /usr/bin | sort | uniq | less
\end{lstlisting}
\end{latin}

در این ساختار، خروجیِ حاصل از \cmd{ls} ابتدا توسط \cmd{sort} به نظم درآمده و سپس \cmd{uniq} نمونه‌های تکراری را از جریان داده حذف می‌کند. شایان ذکر است که \cmd{uniq} برای عملکرد صحیح، حتماً به داده‌های مرتب‌سازی‌شده نیاز دارد؛ چرا که تنها سطرهای تکراریِ مجاور یکدیگر را شناسایی می‌کند.

چنانچه به‌جای حذف تکرارها، صرفاً قصد شناسایی و نمایش موارد مشابه را داشته باشید، می‌توانید از گزینه \cmd{-d} (مخفف \en{duplicate}) استفاده کنید:

\begin{latin}
\begin{lstlisting}
[me@linuxbox ~]$ ls /bin /usr/bin | sort | uniq -d | less
\end{lstlisting}
\end{latin}

این دستور به شما اجازه می‌دهد منحصراً نام برنامه‌هایی را مشاهده کنید که در هر دو دایرکتوری مذکور وجود دارند.

\section{دستور wc: شمارش خطوط، واژگان و بایت‌ها}

فرمان \cmd{wc} که مخفف عبارت \en{Word Count} است، جهت محاسبه‌ی تعداد خطوط، واژگان و حجم داده‌ها (بر حسب بایت) در یک فایل به‌کار می‌رود. ساختار خروجی این دستور به شرح زیر است:

\begin{latin}
\begin{lstlisting}
[me@linuxbox ~]$ wc ls-output.txt
  7902   64566  503634 ls-output.txt
\end{lstlisting}
\end{latin}

در این گزارش، به ترتیب سه شاخص کلیدی ارائه شده است: تعداد خطوط، تعداد کلمات و در نهایت تعداد بایت‌های تشکیل‌دهنده‌ی فایل \file{ls-output.txt}.

مشابه سایر ابزارهای سیستمی، \cmd{wc} در غیاب آرگومانِ فایل، داده‌های مورد نیاز خود را از ورودی استاندارد (\en{stdin}) دریافت می‌کند. با استفاده از گزینه \cmd{-l} (مخفف \en{line})، می‌توان خروجی را منحصراً به شمارش خطوط محدود کرد. این قابلیت در محیط پایپ‌لاین (\en{Pipeline}) بسیار کاربردی است؛ برای نمونه، جهت استخراج آمار دقیق برنامه‌های اجرایی منحصربه‌فرد در مسیرهای \file{/bin} و \file{/usr/bin} از ترکیب زیر استفاده می‌شود:

\begin{latin}
\begin{lstlisting}
[me@linuxbox ~]$ ls /bin /usr/bin | sort | uniq | wc -l
2728
\end{lstlisting}
\end{latin}

در این زنجیره‌ی عملیاتی، ابتدا فهرست فایل‌ها تولید و مرتب‌سازی شده، سپس مدخل‌های تکراری حذف می‌گردند و در نهایت، دستور \cmd{wc -l} با شمارش سطرهای باقی‌مانده، تعداد دقیق فایل‌های متمایز را در خروجی چاپ می‌کند.

\section{دستور grep: جستجو و تطبیق الگوهای متنی}

فرمان \cmd{grep} یکی از قدرتمندترین ابزارهای خط فرمان جهت جست‌وجو و تطبیق الگوهای متنی (\en{Pattern Matching}) در فایل‌ها یا جریان‌های داده است. این دستور، ورودی را به‌صورت سطر به سطر پیمایش کرده و منحصراً خطوطی را که با الگوی تعریف‌شده مطابقت دارند، به خروجی استاندارد (\en{stdout}) ارسال می‌کند.

ساختار کلی این فرمان به شرح زیر است:

\begin{latin}
\begin{lstlisting}
grep pattern [file...]
\end{lstlisting}
\end{latin}

هنگامی که \cmd{grep} در متن ورودی با الگوی مشخص‌شده مواجه می‌شود، کل آن سطر را چاپ می‌کند. الگوهایی که این ابزار قادر به شناسایی آن‌هاست، طیف وسیعی از نویسه‌های ساده تا ساختارهای بسیار پیچیده را شامل می‌شوند. در این مرحله، تمرکز ما بر روی الگوهای متنی ساده (\en{Plain Text}) خواهد بود. الگوهای پیشرفته‌تر که تحت عنوان عبارات باقاعده (\en{Regular Expressions}) شناخته می‌شوند، در مباحث آتی (فصل ۱۹) به‌طور جامع کالبدشکافی خواهند شد.

به عنوان یک نمونه‌ی عملیاتی، فرض کنید قصد داریم تمامی برنامه‌های اجرایی را که نام آن‌ها حاوی عبارت \en{``zip''} است، شناسایی کنیم. این جست‌وجو می‌تواند دیدگاه جامعی درباره‌ی ابزارهای مدیریت فشرده‌سازی موجود در سیستم به ما ارائه دهد. برای این منظور، فرمان \cmd{grep} را به انتهای پایپ‌لاین خود اضافه می‌کنیم:

\begin{latin}
\begin{lstlisting}
[me@linuxbox ~]$ ls /bin /usr/bin | sort | uniq | grep zip
bunzip2
bzip2
gunzip
gzip
unzip
zip
zipcloak
zipgrep
zipinfo
zipnote
zipsplit
\end{lstlisting}
\end{latin}

خروجی شامل تمامی فایل‌هایی خواهد بود که رشته‌ی متنی \en{``zip''} در نام آن‌ها درج شده است. در ادامه، برخی از مهم‌ترین گزینه‌های فنی جهت بهینه‌سازی رفتار \cmd{grep} آورده شده است:

\begin{itemize}
  \item \cmd{-i} (\en{Ignore Case}): فرآیند جست‌وجو را نسبت به بزرگی و کوچکی حروف غیرحساس می‌کند (در حالت پیش‌فرض، \cmd{grep} میان حروف بزرگ و کوچک تمایز قائل می‌شود).
  \item \cmd{-l} (\en{Files with Matches}): به‌جای چاپ خطوط منطبق، صرفاً نام فایل‌هایی را که حداقل یک مورد انطباق در آن‌ها یافت شده است، فهرست می‌کند.
  \item \cmd{-v} (\en{Invert Match}): خروجی را معکوس کرده و تنها خطوطی را نمایش می‌دهد که با الگوی مورد نظر تطابق ندارند.
  \item \cmd{-w} (\en{Word Regexp}): تطبیق را منحصراً به کلمات کامل (\en{Whole Words}) محدود می‌کند تا از شناسایی الگوهایی که بخشی از یک واژه‌ی بزرگتر هستند، جلوگیری شود.
\end{itemize}

\section{دستورهای head و tail: نمایش سطرهای ابتدایی و انتهایی فایل‌ها}

در بسیاری از سناریوها، بررسی تمام محتوای یک فایل ضرورتی ندارد و صرفاً مشاهده‌ی چند سطر نخست یا پایانی کفایت می‌کند. در چنین مواردی، فرامین \cmd{head} و \cmd{tail} ابزارهای بسیار کارآمدی هستند. این دو ابزار به‌طور پیش‌فرض ۱۰ سطر اول یا آخر داده‌ها را نمایش می‌دهند، اما می‌توان این تعداد را با استفاده از گزینه \cmd{-n} شخصی‌سازی کرد.

در مثال زیر، با استفاده از \cmd{head -n 5} تنها پنج سطر ابتدایی فایل \file{ls-output.txt} فراخوانی می‌شود و به همین ترتیب، دستور \cmd{tail -n 5} پنج سطر پایانی فایل مذکور را در خروجی چاپ می‌کند:

\begin{latin}
\begin{lstlisting}
[me@linuxbox ~]$ head -n 5 ls-output.txt
total 343496
-rwxr-xr-x 1 root root      31316 2007-12-05 08:58 [
-rwxr-xr-x 1 root root       8240 2007-12-09 13:39 411toppm
-rwxr-xr-x 1 root root     111276 2007-11-26 14:27 a2p
-rwxr-xr-x 1 root root      25368 2006-10-06 20:16 a52dec
[me@linuxbox ~]$ tail -n 5 ls-output.txt
-rwxr-xr-x 1 root root       5234 2007-06-27 10:56 znew
-rwxr-xr-x 1 root root        691 2005-09-10 04:21 zonetab2pot.py
-rw-r--r-- 1 root root        930 2007-11-01 12:23 zonetab2pot.pyc
-rw-r--r-- 1 root root        930 2007-11-01 12:23 zonetab2pot.pyo
lrwxrwxrwx 1 root root          6 2016-01-31 05:22 zsoelim -> soelim
\end{lstlisting}
\end{latin}

این ابزارها به‌ویژه هنگام کار با پایپ‌لاین‌ها (\en{Pipelines}) برای محدود کردن حجم خروجی فرامین بسیار مفید واقع می‌شوند.

این دستورات در محیط پایپ‌لاین پتانسیل عملیاتی بسیار بالایی دارند و امکان فیلتر کردن دقیق خروجی‌ها را فراهم می‌کنند. برای نمونه، جهت پایش صرفاً پنج سطر انتهاییِ خروجی یک دستور، می‌توان جریان داده را به \cmd{tail -n 5} هدایت کرد:

\begin{latin}
\begin{lstlisting}
[me@linuxbox ~]$ ls /usr/bin | tail -n 5
znew
zonetab2pot.py
zonetab2pot.pyc
zonetab2pot.pyo
zsoelim
\end{lstlisting}
\end{latin}

در این ساختار، پوسته ابتدا خروجی کامل دستور \cmd{ls} را تولید کرده و سپس آن را به عنوان ورودی به \cmd{tail} تحویل می‌دهد. ابزار \cmd{tail} نیز با پردازش این جریان، تمامی سطرها را نادیده گرفته و تنها پنج مورد آخر را در خروجی استاندارد چاپ می‌کند. این رویکرد به‌ویژه زمانی که با فهرست‌های بسیار طولانی سروکار دارید و تنها به نتایج نهایی علاقه‌مند هستید، بسیار کارآمد است.

با ترکیب هوشمندانه ابزارهای \cmd{head} و \cmd{tail} می‌توان بخش‌های میانی یک فایل را با دقت بالایی استخراج کرد. به عنوان مثال، جهت پایش سطرهای ۱۱ تا ۲۰ یک فایل، ابتدا با استفاده از \cmd{head -n 20} بیست سطر نخست را فراخوانی کرده و سپس خروجی حاصل را به \cmd{tail -n 10} هدایت می‌کنیم تا ده سطر پایانی آن مجموعه (که دقیقاً معادل سطر‌های ۱۱ تا ۲۰ فایل اصلی است) نمایش داده شود. همچنین جهت حذف سرآیند‌ها (\en{Headers}) و پانویس‌ها (\en{Footers}) از یک فایل متنی، می‌توان ترکیبی از این دو فرمان را به شکل زیر به کار گرفت. فرض کنید فایلی تحت عنوان \file{text\_header\_footer.txt} در اختیار داریم که حاوی ۵ سطر سرآیند و ۵ سطر پانویس است و قصد داریم تنها بدنه اصلی متن را حفظ کنیم:

\begin{latin}
\begin{lstlisting}
[me@linuxbox ~]$ head -n -5 text_header_footer.txt | tail -n +5 > text.txt
\end{lstlisting}
\end{latin}

در این زنجیره فرامین:

\begin{itemize}
  \item \cmd{head -n -5}: این دستور تمامی سطرهای فایل را به استثنای ۵ سطر انتهایی به خروجی استاندارد ارسال می‌کند.
  \item \cmd{|} (\en{Pipe}): خروجیِ فیلتر شده توسط \cmd{head} را به عنوان ورودی در اختیار \cmd{tail} قرار می‌دهد.
  \item \cmd{tail -n +5}: این دستور از ابتدای ورودی خود، ۵ سطر نخست (سرآیند‌ها) را نادیده گرفته و پردازش را از سطر ششم تا انتها ادامه می‌دهد.
  \item \cmd{>}: در نهایت، برآیند این پردازش‌ها به فایل جدیدی به نام \file{text.txt} بازهدایت (\en{Redirect}) می‌شود.
\end{itemize}

با بهره‌گیری از این متدولوژی، تنها داده‌های اصلی که در میان سرآیند و پانویس محصور شده‌اند، در فایل مقصد ذخیره می‌گردند.

فرمان \cmd{tail} مجهز به گزینه بسیار کاربردی \cmd{-f} (مخفف \en{follow}) است که امکان نظارت بی‌درنگ و زنده بر محتوای یک فایل را فراهم می‌کند. این قابلیت برای مدیران سیستم جهت پایش فایل‌های گزارش (\en{Log Files}) که به‌طور مستمر در حال به‌روزرسانی هستند، ابزاری حیاتی محسوب می‌شود.

به‌عنوان نمونه، برای رصد وقایع سیستم در لحظه، می‌توانید فایل‌های گزارشی نظیر \file{/var/log/messages} یا \file{/var/log/syslog} را به شرح زیر فراخوانی کنید. شایان ذکر است که دسترسی به این مسیرهای حساس سیستمی معمولاً نیازمند برخورداری از امتیازات ابرکاربر (\en{Root}) است:

\begin{latin}
\begin{lstlisting}
[me@linuxbox ~]$ tail -f /var/log/messages
Feb  8 13:40:05 twin4 dhclient: DHCPACK from 192.168.1.1
Feb  8 13:40:05 twin4 dhclient: bound to 192.168.1.4 -- renewal in
1652 seconds.
Feb  8 13:55:32 twin4 mountd[3953]: /var/NFSv4/musicbox exported to
both 192.168.1.0/24 and twin7.localdomain in
192.168.1.0/24,twin7.localdomain
Feb  8 14:07:37 twin4 dhclient: DHCPREQUEST on eth0 to 192.168.1.1
port 67
Feb  8 14:07:37 twin4 dhclient: DHCPACK from 192.168.1.1
Feb  8 14:07:37 twin4 dhclient: bound to 192.168.1.4 -- renewal in
1771 seconds.
Feb  8 14:09:56 twin4 smartd[3468]: Device: /dev/hda, SMART Prefailure Attribute: 8 Seek_Time_Performance changed from 237 to 236
Feb  8 14:10:37 twin4 mountd[3953]: /var/NFSv4/musicbox exported to both 192.168.1.0/24 and twin7.localdomain
192.168.1.0/24,twin7.localdomain
Feb  8 14:25:07 twin4 sshd(pam_unix)[29234]: session opened for user me by (uid=0)
Feb  8 14:25:36 twin4 su(pam_unix)[29279]: session opened for user root by me(uid=500)
\end{lstlisting}
\end{latin}

با اجرای این دستور، \cmd{tail} ابتدا ۱۰ سطر پایانی فایل را چاپ کرده و سپس در حالت تعلیق باقی می‌ماند. به محض اضافه شدن ورودی‌های جدید توسط سیستم به انتهای فایل، \cmd{tail} بلافاصله آن‌ها را در خروجی استاندارد نمایش می‌دهد. این فرآیند پایش تا زمان توقف دستی توسط کاربر با استفاده از ترکیب کلیدی \cmd{Ctrl+c} ادامه خواهد یافت.

\section{دستور tee: انشعاب جریان داده در پایپ‌لاین‌ها}

در اکوسیستم خط فرمان لینوکس، ابزار \cmd{tee} نقشی مشابه یک اتصال \en{``T''} شکل در سیستم‌های پایپ‌لاین ایفا می‌کند. این برنامه به‌طور هم‌زمان داده‌ها را از ورودی استاندارد (\en{stdin}) دریافت کرده و یک نسخه از آن را به خروجی استاندارد (\en{stdout}) جهت استمرار در پایپ‌لاین و نسخه‌ای دیگر را در یک یا چند فایل مشخص کپی می‌کند. این قابلیت برای مستندسازی یا ذخیره‌سازی داده‌ها در مراحل میانی یک فرآیند پردازشی بسیار ارزشمند است.

برای تبیین این عملکرد، یکی از مثال‌های پیشین را با الحاق دستور \cmd{tee} بازنگری می‌کنیم. در این سناریو، قصد داریم فهرست کامل فایل‌ها را در فایل \file{ls.txt} ذخیره کنیم و هم‌زمان، آن جریان را برای فیلتر شدن توسط \cmd{grep} در مسیر پایپ‌لاین حفظ نماییم:

\begin{latin}
\begin{lstlisting}
[me@linuxbox ~]$ ls /usr/bin | tee ls.txt | grep zip
bunzip2
bzip2
gunzip
gzip
unzip
zip
zipcloak
zipgrep
zipinfo
zipnote
zipsplit
\end{lstlisting}
\end{latin}

در این زنجیره فرامین:

\begin{itemize}
  \item \cmd{ls /usr/bin}: محتویات دایرکتوری را استخراج کرده و به جریان خروجی می‌فرستد.
  \item \cmd{tee ls.txt |}: داده‌ها را از مرحله قبل دریافت نموده، ضمن ذخیره در فایل \file{ls.txt} (بدون وقفه در جریان)، آن‌ها را به مرحله بعد هدایت می‌کند.
  \item \cmd{grep zip |}: ورودیِ دریافتی از \cmd{tee} را برای یافتن الگوی \en{``zip''} پالایش کرده و نتایج نهایی را در ترمینال چاپ می‌کند.
\end{itemize}

این مثال به‌روشنی نشان می‌دهد که \cmd{tee} چگونه یک جریان داده را بدون ایجاد تغییر در مسیر اصلی، به یک مقصد جانبی (فایل) نیز منشعب می‌کند. مشابه سایر ابزارهای بازهدایت، \cmd{tee} نیز با استفاده از گزینه \cmd{-a} (مخفف \en{append}) قادر است داده‌ها را به‌جای بازنویسی (\en{Overwrite})، به انتهای فایل موجود الحاق نماید.

\section{جمع‌بندی}

همان‌طور که همواره تأکید شده است، مطالعه مستندات رسمی (\en{Manual Pages}) برای هر یک از فرامینی که در این فصل بررسی کردیم، امری ضروری است. آنچه بازگو شد، تنها سطح ابتدایی کاربرد این ابزارها بود؛ در حالی که هر یک از آن‌ها از گزینه‌ها و قابلیت‌های پیشرفته‌ای برخوردارند که با افزایش تجربه عملی، به‌تدریج بر آن‌ها مسلط خواهید شد.

مکانیزم بازهدایت (\en{Redirection})، ابزاری بنیادین و بی‌بدیل برای حل چالش‌های فنی در محیط خط فرمان محسوب می‌شود. اکثر فرامین سیستمی از ورودی و خروجی استاندارد بهره می‌گیرند و تقریباً تمامی برنامه‌های خط فرمان، پیام‌های وضعیتی و گزارش‌های خطا را از طریق مجرای خطای استاندارد (\en{stderr}) منتشر می‌کنند.

\section{تحلیل تطبیقی لینوکس و ویندوز: از مصرف‌کننده تا معمار}

برای درک بهتر تفاوت بنیادین میان فلسفهٔ ویندوز و لینوکس، بیایید از یک قیاس ساده در دنیای سرگرمی کمک بگیریم.

ویندوز را مانند یک کنسول بازی اختصاصی تصور کنید؛ مثلاً یک \en{Game Boy}. شما دستگاهی زیبا و آماده‌به‌کار می‌خرید که با رابط گرافیکی جذاب و جلوه‌های صوتی از پیش‌طراحی‌شده، تجربه‌ای لذت‌بخش را در اختیارتان می‌گذارد. اما رفته‌رفته متوجه می‌شوید که تنها به بازی‌هایی که برای این دستگاه ساخته شده‌اند، محدود هستید و برای هر قابلیت جدید ناچارید به سازنده بازگردید. این چرخه آنقدر ادامه می‌یابد تا روزی به دنبال کاری خاص باشید و بشنوید که «چنین قابلیتی ممکن نیست، چون بازار خواهان آن نبوده است». در این مدل، شما تنها مصرف‌کننده‌ای هستید که با قواعد از پیش‌تعیین‌شده بازی می‌کنید.

در سوی دیگر، لینوکس مانند بزرگ‌ترین جعبه‌ابزار مهندسی جهان است؛ همانند مجموعه‌های بی‌نهایت لِگو. در ابتدا با صدها قطعهٔ جداگانه، چرخ‌دنده، موتور و نقشه‌های راهنما روبه‌رو می‌شوید. کم‌کم با کنار هم چیدن این قطعات، نه تنها می‌توانید سازه‌های پیشنهادی را بسازید، بلکه خلاقیت خود را به کار گرفته و هر آنچه در ذهن دارید، طراحی و پیاده‌سازی می‌کنید. در این فضا، هیچ وابستگی به تأمین‌کنندهٔ بیرونی ندارید؛ تمام ابزارهایی که برای خلق ایده‌های فنی‌تان نیاز دارید، در اختیار شماست. این سیستم دقیقاً به شکلی درمی‌آید که شما آن را مهندسی کرده‌اید.

انتخاب میان این دو، به نیاز و سلیقهٔ شخصی شما بستگی دارد. اما پرسش اصلی این است: کدام تجربه برای شما ارزشمندتر است؟ لذت بردن از محصولی کامل و از پیش‌ساخته، یا داشتن ابزاری بی‌نهایت که قدرت آفرینش را به دست شما می‌دهد؟